\documentclass[10pt,fleqn,a4paper]{jsarticle}
\usepackage[dvipdfmx]{graphicx,color}
\usepackage{ascmac,amsmath,amssymb,amstext}
\usepackage{tikz,multicol,float,tikz-3dplot}
\usetikzlibrary{positioning,intersections,calc,arrows.meta,math,angles}
\usepackage{zogeny,ceo,comment}
\newcommand{\ans}[1]{\mbox{\boldmath{$#1$}}}
\renewcommand{\baselinestretch}{1.2} 
\tdplotsetmaincoords{60}{110}%{xy平面からどれだけ上にいるか(90度から-90度)}{z軸中心にどれだけ左右に回転するか}
\title{特殊漸化式の基本}
\author{大島　遙斗}
\begin{document}
\maketitle
\newpage
基本漸化式をこの間扱いました.今回は,特殊な漸化式の発想・解法を扱います.\ans{特殊}と銘打っていますが,
解き方は基本的な3つの漸化式に帰着させることです.その帰着の方法が,\ans{特殊}なだけです.\\
それでは,順にみてゆきましょう.それでは,参ります.\\\\
$\bullet\doichi\ans{和で表された漸化式.}$\\
次のような形をした漸化式です.
\begin{center}
    $S_n\p{=\wa{k=p}{q}a_k}=ma_n+f(n)$
\end{center}
\kangaekata\\
数列の構造を確かめれば簡単です.次のような構造になっています.
\begin{center}
    $S_n=a_1+a_2+a_3+\cdots+a_n　(n\geq1)$
\end{center}
ここで,上の構造から$a_nのみ残したかったら,\ans{a_1\sim a_{n-1}を消せばよい!}ということがわかります.したがって,S_{n-1}を考えましょう.$
\begin{align*}
    S_n&=a_1+a_2+\cdots+a_{n-1}+a_n　(n\geq1)\cdots\cdots\cdots\maruichi\\
    S_{n-1}&=a_1+a_2+\cdots+a_{n-1}\hspace{13mm}(n\geq2)\cdots\cdots\cdots\maruni
\end{align*}
$\maruichi-\maruni を考えると,a_nが残るので,n\geq2 の範囲で,a_n のみの漸化式が得られる.$\\
では,実際に次の2問を10分を目安に解きましょう.\\\\
\ba{1}.次の漸化式で定まる数列$\B{a_n}の一般項を求めよ.$\\
$\kakkoichi S_n=-2a_n-2n+5　(n\geq1)$\\
$\kakkoni \wa{k=1}{n}a_k=2a_n+n　(n\geq1)$
\newpage
\noindent\ba{1}.\\
\noindent\kai\\
$\kakkoichib　 S_{n-1}=-2a_{n-1}-2n+7より,S_{n}-S_{n-1}を考えると,S_{n}-S_{n-1}=a_nから$
\begin{align*}
    a_n=-2a_n+2a_{n-1}-2&\doti 3a_n=2a_{n-1}-2\cdots\cdots\cdots\cdots\cdots\cdots\cdots\cdot\maruichi\\
    &\doti a_n=\dfrac{2}{3}a_n-\dfrac{2}{3}　(n\geq2)\cdots\cdots\cdots\maruni 
\end{align*}
このことより,
\begin{align*}
    \maruni&\doti a_n+2=\dfrac{2}{3}\p{a_{n-1}+2}　(n\geq2)\\
    &\doti a_{n+1}+2=\dfrac{2}{3}\p{a_{n}+2}　(n\geq1)\chu インデントをズラしても一般性を失いません.\\
    &\doti a_n+2=\p{\dfrac{2}{3}}^{n-1}\cdot\p{a_1+2}=\p{\dfrac{2}{3}}^{n-1}\cdot3\\
    &\doti \ans{a_n=3\cdot\p{\dfrac{2}{3}}^{n-1}-2　(n\geq1)}
\end{align*}\\
\kakkonib　$\wa{k=1}{n}a_k=S_nとおくと,S_{n-1}を考えて,S_n-S_{n-1}=a_nより,$
\begin{align*}
    a_n=2a_n-2a_{n-1}+1&\doti a_n=2a_{n-1}-1　(n\geq2)\\
    &\doti a_{n+1}=2a_n-1　(n\geq1)\\
    &\doti a_{n+1}-1=2\p{a_n-1}\\
    &\doti a_n=\p{a_1-1}\cdot2^{n-1}+1\\
    &\doti a_n=-2\cdot2^{n-1}+1\\
    &\doti \ans{a_n=-2^n+1　(n\geq1)}
\end{align*}
\newpage
\noindent$\bullet\doni\ans{冪乗,積が含む漸化式.}$\\
次のような形をした漸化式です.
\begin{align*}
    &\bullet\hspace{2mm}\ans{a_{n+1}=a_n^k}\cdots\cdots\cdots\cdots\cdots\cdots\maruichi\\\\
    &\bullet\hspace{2mm}\ans{a_na_{n+1}=k\sqrt[p]{a_n}}\cdots\cdots\cdots\maruni
\end{align*}
ここらへんの漸化式は,\ans{嫌なものを取り除く.}という方針で解きます.順にみてゆきましょう.\\
\kangaekata\\
$\maruichi:右辺の\ans{k乗が嫌!}です.指数を前に出せるものといえば....そう,\ans{対数}です.\\では,底は何にしましょうか.まあなんでも良いのですが,\ans{初項を底にする.}と
楽なことが多いです.\\
\maruni:まあこれも上と同じですね.両辺に対数をとりましょう.この場合,\ans{底を冪乗(p)}にするのが楽になる場合が多いです.$\\
では,次の2題を20分を目安に取り組んでみてください.\\\\
\ba{2$\bullet$1}次の条件によって定められる数列$\B{a_n}がある.$
\begin{align*}
    \begin{cases}
        a_1=2\\
        a_{n+1}=8a_n^2　(n=1,2,3\cdots)\cdots\cdots\cdots\cdots\maruichi
    \end{cases}
\end{align*}
\kakkoichi 　$数列\B{a_n}の一般項を求めよ.$\\
\kakkoni　$\mathrm{P}_n=a_1a_2a_3\cdots a_nとおく.数列\B{\mathrm{P}_n}の一般項を求めよ.$\\
\hspace{10cm}(大阪大)\\\\
\ba{2$\bullet$2}次の漸化式で定義される数列$\B{a_n}$の一般項$a_n$を求めよ.
\begin{align*}
    \begin{cases}
        a_1=2\\
        a_na_{n+1}=2\sqrt{a_n}　(n\geq1)\cdots\cdots\cdots\bigstar
    \end{cases}
\end{align*}
\newpage
\noindent\ba{2$\bullet$1}\\
\kangaekata\\
\kakkoichi 両辺に底を2とする対数を取りましょう.\\
\kakkoni 積の形は嫌なので,両辺に任意の底を持つ対数を取りましょう.(解答では,底2を取ります.)\\
\kai\\
\kakkoichib　$\maruichi の両辺に底を2とする対数をとると,$
\begin{align*}
    \log_2a_{n+1}&=\log_2a_n+\log_28\\
    &=2\log_2a_n+3\cdots\cdots\cdots\cdots\maruni
\end{align*}
これより,$b_n=\log_2a_nとおくと,$
\begin{align*}
    \maruni&\doti b_{n+1}=2b_n+3\\
    &\doti b_{n+1}+3=2\p{b_n+3}\\
    &\doti b_n=\p{b_1+3}\cdot2^{n-1}-3\\
    &\doti \log_2a_n=\p{\log_22+3}-3=2^{n+1}-3\\
    &\doti \ans{a_n=2^{2^{n+1}-3}　(n\geq1)}
\end{align*}\\
\kakkonib　$\mathrm{P}_n=a_1a_2a_3\cdots a_n\cdots\cdots\cdots\marusan の両辺に底を2とする対数をとると,$
\begin{align*}
    \log_2\mathrm{P}_n&=\log_2 a_1+\log_2 a_2+\log_2 a_3+\cdots+\log_2 a_n\\
    &=\wa{k=1}{n}\log_2 a_k\\
    &=\wa{k=1}{n}\log_2\p{2^{2^{k+1}-3}}\\
    &=\wa{k=1}{n}\p{2^{k+1}-3}\log_22\\
    &=\dfrac{2^{n+2}-4}{2-1}-3n\\
    &=2^{n+2}-4-3n
\end{align*}
$\yueni \ans{\mathrm{P}_n=2^{2^{n+2}-3n-4}　(n\geq1)}$
\newpage
\noindent\ba{2$\bullet$2}\\
\kangaekata\\
 もちろん積の形が嫌なので,両辺に対数を取ります.底は,漸化式の係数と平方根から2が良いでしょう.\\
 \kai\\
 $\bigstar の両辺に底を2とする対数をとると,$
 \begin{align*}
    \log_2a_n+\log_2a_{n+1}=\log_22\sqrt{a_n}=1+\dfrac{1}{2}\log_2a_n
 \end{align*}
 $\yueni \log_2a_{n+1}=-\dfrac{1}{2}\log_2a_n+1\cdots\cdots\cdots\heartsuit$\\
 $これより,\heartsuit を解くと,$
 \begin{align*}
    \heartsuit&\doti\log_2a_{n+1}-\dfrac{2}{3}=-\dfrac{1}{2}\p{\log_2a_n-\dfrac{2}{3}}\\
    &\doti \log_2a_n=\dfrac{1}{2}\cdot\p{-\dfrac{1}{2}}^{n-1}+\dfrac{2}{3}\\
    &\doti \ans{a_n=2^{\frac{1}{3}\B{\p{-\frac{1}{2}}^{n-1}+2}}　(n\geq1)}
 \end{align*}
 \newpage
\noindent$\bullet\dosan\ans{分数型漸化式.}$\\
名の通り,分数の形をした漸化式です.次のような形.
\begin{center}
    $\ans{a_{n+1}=\dfrac{pa_n+f(n)}{qa_n+f(n)}}$
\end{center}
分数型の漸化式は,次のような解き方があります.前者は,入試問題の典型的な解法です.\\
\kuromaruichi 「$nに順に数を代入して,一般項を推測して数学的帰納法で正当化する.$」\\
\kuromaruni 「$式を試行錯誤して,自分が解ける形に変形する.$」\\
例題を上の二つの方法で解いてみましょう.\\\\
\reidaic　次の漸化式で表される数列$\B{a_n}の一般項a_nを求めよ.$
\begin{align*}
    \begin{cases}
        a_1=2\\
        a_{n+1}=\dfrac{a_n}{a_n+4}　(n\geq1)
    \end{cases}
\end{align*}
\end{document}